\section{Introduction}
\label{sec:intro}

Web agents --- language models that drive a real browser through a ReAct-style
loop of reasoning, tool calls, and page observations --- are now benchmarked
almost exclusively by one number: the fraction of tasks whose final environment
state satisfies the benchmark's functional evaluator
\citep{zhou2024webarena, koh2024visualwebarena}. That number is then compared
across models, across prompting strategies, and --- increasingly --- across
fault-injection conditions, on the implicit assumption that a small drop in
success rate means the agent is robust and a large drop means it is fragile.

This paper asks whether that assumption holds. We take a single agent
configuration (a LangGraph ReAct agent, \ToolCount{} WebArena tools, an
accessibility-tree observation channel) on the WebArena shopping, Reddit, and
GitLab sites, and subject it to a paired fault-injection experiment: for each of
five non-adversarial faults, every (task, seed) cell is run twice, once clean and
once with the fault injected at a fixed step, and both runs are scored by the
official evaluator. The pairing is what makes the design informative --- it
removes task difficulty and per-seed variance from the comparison, and it lets us
run an exact McNemar test rather than comparing two independent proportions.

The headline result is negative, and that is the point. After Holm correction,
\emph{no fault has a detectable effect on official success}
($p_{\text{Holm}}=1.000$ for all five). The single largest point estimate,
\ResAgentParamErrorDeltaNat{}~pp for tool parameter errors, is comfortably inside a
bootstrap interval that spans zero. Read the usual way, this would be reported as
``the agent is robust to these faults.'' We show that reading is not supported:
with only \ResWebDomMissingNNat--\ResWebTimeoutNNat{} native-evaluator pairs per
fault, the design has essentially no power, and the honest statement is that the
experiment could not have detected a moderate effect had one existed.

Where the faults \emph{are} visible is one level below the task outcome. In a
\TotalTrials{}-trial behavioural record that trades the evaluator for step and
wall-clock instrumentation, HTTP errors add \ResHttpErrorStepDelta{} ReAct steps
and \ResHttpErrorTimeDelta{}~s per pair, and tool parameter errors add
\ResParamErrStepDelta{} steps --- both far beyond noise. End-to-end success is
blind to these effects; process-level measurement is not. This is the paper's
central methodological claim: for web agents, task success is a blunt and
high-variance instrument, and robustness conclusions drawn from it alone are
under-determined.

Two further results shape how such experiments should be run. First, the
evaluator path matters more than any fault we inject: a compatibility fallback
handling \OfficialFallbackN{} of the \OfficialN{} official evaluations fails by
construction, and pooling rather than stratifying those trials moves the frozen
baseline's success rate from \BaseNativeRate{} to \BasePooledRate{} --- a
\BaselineSwing{}~pp swing driven purely by instrumentation. Second, we document
degeneracies in the auxiliary artifacts (a recovery table whose columns are
constant or defined by the outcome it claims to explain; behavioural counts whose
``\TotalTrials{} trials'' are in fact \MatrixPairs{} pairs), which are easy to
introduce and, once in a pipeline, easy to read as evidence.

We make three contributions:
\begin{itemize}
  \item \textbf{A paired, evaluator-stratified attribution protocol} for web-agent
  fault injection, with exact and mid-$p$ McNemar tests, pair-unit bootstrap
  intervals, Wilcoxon/Hodges--Lehmann tests on process deltas, and explicit
  multiplicity control (\secref{sec:method}).
  \item \textbf{An empirical demonstration that the standard outcome measure is
  underpowered} at realistic sample sizes, quantified with minimum detectable
  effects and the sample sizes that would actually be needed
  (\secref{sec:results}).
  \item \textbf{Evidence that the measurement instrument can dominate the effect
  being measured}, together with a stratified analysis that removes the
  confound, and a staged design that would make the model-versus-architecture
  question answerable (\secref{sec:discussion}).
\end{itemize}

We are explicit about scope: the data cover one model and one architecture, so we
make \emph{no} claim about how much of the robustness gap is attributable to the
model versus the architecture, which was the motivating question.
\secref{sec:discussion} states what the present design can and cannot identify.
